\begin{table}[t]
    \centering
    \resizebox{0.9\textwidth}{!}{%
    \begin{tabular}{@{}lccccccc@{}}
        \toprule
                           &                 & \multicolumn{2}{c}{\xnli} & \multicolumn{2}{c}{NER} & \multicolumn{2}{c}{\sib} \\
        \cmidrule(lr){3-4} \cmidrule(lr){5-6} \cmidrule(lr){7-8}
        Max subword length & Fertility ratio & \mbert & \xlmr & \mbert & \xlmr & \mbert & \xlmr \\
        \midrule
        6 chars & 1.1--1.2 & 2.8  & 1.8  & 5.2  & 3.0  & 2.9  & 5.4 \\
        4 chars & 1.2--1.5 & 10.2 & 9.9  & 14.3 & 10.6 & 18.6 & 23.6 \\
        3 chars & 1.4--1.7 & 16.8 & 18.3 & 24.4 & 19.3 & 42.1 & 40.9 \\
        2 chars & 1.8--2.3 & 30.6 & 30.2 & 39.2 & 33.1 & 54.5 & 59.4 \\
        1 char  & 2.9--3.9 & 40.7 & 44.3 & 58.7 & 56.4 & 63.6 & 67.4 \\
        \bottomrule
    \end{tabular}
    }
    \caption{
        Drop in English test score (points) when we cap the length of every subword at test time.
        The fertility ratio is relative to standard English tokenization and varies by model and task.
        Values are means over seeds; seed SD is under 1 point for \xnli, up to 4 for NER and up to 9 for \sib.
        At the 1-character cap, 8\% (\sib) and 18--19\% (\xnli) of tokens are also truncated by the length limit.
    }
    \label{tab:intervention}
\end{table}
